% Satz des Pythagoras – bank/pythagoras/ (zusammenbau v0.4, Heft)
\einheitenkopf[t4]{Satz des Pythagoras}
\setcounter{aufgabe}{26}

% Anlauf (ohne Prüfkennung)
\begin{aufgabe}{Hypotenuse – Anlauf}
\begin{teile}
\teil Rechten Winkel einkreisen, Hypotenuse mit H – wie lang ist sie?

\dreieckrw{5}{1.12}{$40$ cm}{$9$ cm}{$41$ cm}
\leerfeld[cm]
\teil Katheten $a = 30$ cm, $b = 40$ cm – Hypotenuse $c$? c² = \leerfeld , c = \leerfeld[cm]
\teil Welche Gleichung gilt für $g$? \\ \kreuz{$g = \sqrt{e^2 + f^2}$} \\ \kreuz{$g = e^2 + f^2$} \\ \kreuz{$g = \sqrt{e^2 - f^2}$} \\ \kreuz{$g = e + f$}

\dreieckrw{5}{3.57}{$e$}{$f$}{$g$}
\end{teile}
\end{aufgabe}

\begin{aufgabe}{Hypotenuse – Prüfungsaufgaben}
\begin{teile}
\teil Ein dreieckiges Grundstück ABC hat bei C einen rechten Winkel. Der Zaun von B nach C ist $14$ m lang, der von A nach C $39$ m. Berechne die Länge der Seite AB. Gib das Zwischenergebnis $AB^2$ an und runde auf eine Dezimale \hfill (P10 2020 OS).

\dreieckrw{5}{1.79}{$39$ m}{$14$ m}{?}
AB² = \leerfeld , AB ≈ \leerfeld[m]
\teil Ein Segelboot fährt von der Boje A zur Boje C, dort biegt es im rechten Winkel ab zur Boje B. Es gilt $AC = 47$ m und $CB = 22$ m. Berechne die direkte Entfernung AB. Gib das Zwischenergebnis $AB^2$ an und runde auf eine Dezimale \hfill (P10 2020 OS).

\dreieckrw{5}{2.34}{$47$ m}{$22$ m}{?}
AB² = \leerfeld , AB ≈ \leerfeld[m]
\teil Eine Laderampe führt auf einer waagerechten Länge von $240$ cm um $18$ cm nach oben. Berechne die Länge der Rampe auf eine Dezimale \hfill (P10 2025 OS).

\dreieckrw{6}{0.45}{$240$ cm}{$18$ cm}{?}
\leerfeld[cm]
\teil Ein Steg zu einem Boot überwindet auf $350$ cm waagerechter Länge einen Höhenunterschied von $22$ cm. Berechne die Länge des Stegs auf eine Dezimale \hfill (P10 2025 OS).

\dreieckrw{6}{0.38}{$350$ cm}{$22$ cm}{?}
\leerfeld[cm]
\end{teile}
\end{aufgabe}

% Anlauf (ohne Prüfkennung)
\begin{aufgabe}{Kathete – Anlauf}
\begin{teile}
\teil Leiterlänge und Abstand des Leiterfußes zur Wand sind gegeben, gesucht ist die Höhe an der Wand. Plus oder minus? \\ \kreuz{Hypotenuse gesucht – plus} \\ \kreuz{Kathete gesucht – minus}
\teil Hypotenuse $c = 89$ cm, Kathete $b = 80$ cm – Kathete $a$? a² = \leerfeld , a = \leerfeld[cm]
\teil Welche Gleichung gilt für die Kathete $f$? \\ \kreuz{$f = \sqrt{g^2 - e^2}$} \\ \kreuz{$f = \sqrt{g^2 + e^2}$} \\ \kreuz{$f = g^2 - e^2$} \\ \kreuz{$f = g - e$}

\dreieckrw{5}{3.57}{$e$}{$f$}{$g$}
\end{teile}
\end{aufgabe}

\begin{aufgabe}{Kathete – Prüfungsaufgaben}
\begin{teile}
\teil Eine Seilrutsche ist $426$ m lang und überbrückt waagerecht $390$ m. Berechne den Höhenunterschied zwischen Start und Ziel auf eine Dezimale \hfill (P10 2024 OS).

\dreieckrw{5}{2.2}{$390$ m}{$h$}{$426$ m}
\leerfeld[m]
\teil Die Trasse eines Sessellifts ist $1\,250$ m lang, waagerecht gemessen sind es $1\,180$ m. Berechne, wie viele Meter der Lift nach oben fährt, auf eine Dezimale \hfill (P10 2024 OS).

\dreieckrw{5}{1.75}{$1\,180$ m}{$h$}{$1\,250$ m}
\leerfeld[m]
\teil Im Dreieck PQR ist der Winkel bei R stumpf. Die Höhe von R auf die Seite PQ ist $5{,}3$ m lang, ihr Fußpunkt F liegt auf PQ. Es gilt $PR = 9{,}6$ m. Weise nach, dass $PF \approx 8{,}0$ m ist \hfill (P10 2022 OS).
\teil Im Dreieck UVW ist der Winkel bei W stumpf. Die Höhe von W auf die Seite UV ist $5{,}7$ cm lang, ihr Fußpunkt S liegt auf UV. Es gilt $UW = 13{,}6$ cm. Weise nach, dass $US \approx 12{,}3$ cm ist \hfill (P10 2022 OS).
\steil Ein Fenster hat die Form eines Drachenvierecks ABCD. Seine Diagonalen schneiden sich in M rechtwinklig, und es gilt $AM = MC = DM = 32$ cm. Begründe mit der Umkehrung des Satzes des Pythagoras, dass das Fenster bei D einen rechten Winkel hat \hfill (P10 2025 OS).
\steil Ein dreieckiges Segel PQR: Der Punkt S liegt auf PR, die Strecke QS steht senkrecht auf PR. Es gilt $PS = 1{,}6$ m, $SR = 0{,}9$ m und $QS = 1{,}2$ m. Begründe mit der Umkehrung des Satzes des Pythagoras, dass das Segel bei Q einen rechten Winkel hat \hfill (P10 2025 OS).
\end{teile}
\end{aufgabe}

% Anlauf (ohne Prüfkennung)
\begin{aufgabe}{Figuren und Körper – Anlauf}
\begin{teile}
\teil Raute mit Diagonalen: ein rechtwinkliges Teildreieck nachfahren. Welche Seiten sind Katheten, welche ist die Hypotenuse?

\raute[diagonalen]{4}{2.5}
\teil Gleichschenkliges Dreieck: Grundseite $18$ cm, Schenkel $41$ cm – Höhe $h$?

\dreieck{(0,0)}{(2.02,0)}{(1.01,4.5)}{$41$ cm}{}{$18$ cm}{}{}{}
h² = \leerfeld , h = \leerfeld[cm]
\teil $E(-4 | -3)$ und $F(5 | 9)$ – Länge der Strecke EF? EF = \leerfeld
\end{teile}
\end{aufgabe}

\begin{aufgabe}{Figuren und Körper – Prüfungsaufgaben}
\begin{teile}
\steil Ein Pavillon hat ein kegelförmiges Zeltdach mit dem Durchmesser $7{,}2$ m. Die Höhe des Dachs ist bekannt, steht hier aber nicht. Beschreibe, wie man die Länge einer Zeltstange vom Dachrand bis zur Spitze berechnet \hfill (P10 2018 OS).

\kegel{1.8}{1.5}{}{$h$}{$s$}
\steil Ein Kirchturm ist insgesamt $42$ m hoch und trägt eine kegelförmige Spitze mit dem Durchmesser $5{,}8$ m. Die Höhe der Spitze allein ist bekannt, steht hier aber nicht. Beschreibe, wie man die Länge eines Sparrens von der Traufe bis zur Spitze berechnet \hfill (P10 2018 OS).

\kegel{1.2}{2.5}{}{$h$}{$s$}
\teil Ein Getreidesilo besteht aus einem Zylinder mit der Höhe $19{,}5$ m und einem Kegeldach mit Radius $3{,}9$ m und Mantellinie $6{,}2$ m. Berechne die Gesamthöhe des Silos auf eine Dezimale \hfill (P10 2026 FOR).

\zylinder{1}{3}{}{$19{,}5$ m}
\leerfeld[m]
\teil Ein Leuchtturm besteht aus einem zylindrischen Schaft mit der Höhe $31{,}0$ m und einem Kegeldach mit Radius $2{,}6$ m und Mantellinie $5{,}1$ m. Berechne die Gesamthöhe des Leuchtturms auf eine Dezimale \hfill (P10 2026 FOR).

\zylinder{1}{3}{}{$31$ m}
\leerfeld[m]
\steil Ein zylindrisches Glas ist innen $11$ cm hoch und hat den Radius $2{,}8$ cm. Ein Strohhalm steht schräg von der Bodenkante zur gegenüberliegenden oberen Kante und ragt $4$ cm heraus. Berechne die Länge des Strohhalms auf eine Dezimale \hfill (P10 2022 OS).

\zylinder{0.76}{3}{$2{,}8$ cm}{$11$ cm}
\leerfeld[cm]
\steil Eine zylindrische Vase ist innen $24$ cm hoch und hat den Radius $4{,}5$ cm. Ein Zweig steht schräg von der Bodenkante zur gegenüberliegenden oberen Kante und ragt $6$ cm heraus. Berechne die Länge des Zweigs auf eine Dezimale \hfill (P10 2022 OS).

\zylinder{0.56}{3}{$4{,}5$ cm}{$24$ cm}
\leerfeld[cm]
\end{teile}
\end{aufgabe}

\begin{aufgabe}{Figuren und Körper – Prüfungsaufgaben – weiter}
\begin{teile}
\teil Die Punkte R und S sind im Koordinatensystem eingezeichnet. Berechne die Länge der Strecke RS auf eine Dezimale \hfill (P10 2019 OS).

\begin{ksys}[xmin=-4,xmax=5,ymin=-4,ymax=4,ablesen] \punkt{-2}{3}{R} \punkt{4}{-1}{S} \end{ksys}
RS ≈ \leerfeld
\teil Die Punkte G und J sind im Koordinatensystem eingezeichnet. Berechne die Länge der Strecke GJ auf eine Dezimale \hfill (P10 2019 OS).

\begin{ksys}[xmin=-4,xmax=5,ymin=-4,ymax=4,ablesen] \punkt{1}{-3}{G} \punkt{-3}{2}{J} \end{ksys}
GJ ≈ \leerfeld
\end{teile}
\end{aufgabe}
